\section{Document retrieval on the endpoint: \textit{corpus}}
\label{sec:corpus}

\subsection{What it is}
\textit{corpus} is a retrieval-augmented question answering system over our own
documents, running entirely on American Terawatt hardware and using the
endpoint of Section~\ref{sec:remote-access} as its model. Documents are parsed
once, split into passages, embedded, and indexed on local disk. At question
time the relevant passages are retrieved and handed to \cmd{gpt-oss:120b} on
the shared GPU pool, which answers from those passages and cites them. Nothing
leaves the network, and nothing is sent to an external API.

The system deliberately does \emph{not} fine-tune a model on the documents.
Fine-tuning gives no reliable recall of specific facts, no citations, and would
need re-running whenever a document changes. Retrieval at query time gives all
three.

\subsection{Where everything runs}
Three machines are involved, and the split matters: the corpus lives next to
the documents, the model lives on the GPU, and only questions and answers cross
the network.

\begin{figure}[!ht]
  \centering
  \resizebox{\linewidth}{!}{%
  \begin{tikzpicture}[node distance=9mm,
    blk/.style={draw=atBorderStrong, rounded corners=2pt, fill=atBg,
                align=center, inner sep=4pt, font=\scriptsize,
                minimum height=8mm},
    store/.style={blk, fill=atBgSubtle},
    comp/.style={blk, fill=atBgMuted},
    model/.style={blk, fill=atAccent!30},
    zone/.style={draw=atBorderStrong, dashed, rounded corners=4pt, inner sep=6mm},
    flow/.style={-{Stealth[length=2mm]}, draw=atFgMuted, thick},
    net/.style={-{Stealth[length=2mm]}, draw=atFgMuted, thick, dashed},
    lbl/.style={font=\tiny, text=atFgMuted, align=center},
  ]
  % laptop
  \node[comp] (hermes) {Hermes Agent\\\textnormal{(or any MCP client)}};
  \node[zone, fit=(hermes), label={[font=\small\bfseries]above:Your laptop}] (zlap) {};

  % developer VM
  \node[comp, right=30mm of hermes] (server) {corpus server\\\cmd{:8765}\\MCP + JSON};
  \node[comp, above right=6mm and 12mm of server] (cli) {\cmd{./rag} CLI};
  \node[comp, below right=6mm and 12mm of server] (core)
    {\textbf{ragkit}\\parse $\cdot$ chunk\\retrieve $\cdot$ answer};
  \node[store, below=of core] (lance) {LanceDB index\\vectors + BM25};
  \node[store, below=of lance] (md) {\cmd{md/} parse cache\\keyed on sha256};
  \node[store, left=9mm of md] (raw) {\cmd{\char`~/work/raw}\\dashboard uploads};
  \node[zone, fit=(server)(cli)(core)(lance)(md)(raw),
        label={[font=\small\bfseries]above:dev-eugen \textnormal{(developer VM)}}] (zdev) {};

  % GPU host
  \node[comp, right=30mm of core] (ollama) {Ollama\\\cmd{:11434}};
  \node[model, above right=4mm and 10mm of ollama] (chat)
    {\cmd{gpt-oss:120b}\\65.4\,GB MXFP4\\131k ctx, tools};
  \node[model, below right=4mm and 10mm of ollama] (embed) {\cmd{bge-m3}\\1024-dim};
  \node[zone, fit=(ollama)(chat)(embed),
        label={[font=\small\bfseries]above:GPU host \textnormal{(pool container)}}] (zgpu) {};

  % edges
  \draw[net] (hermes.east) -- node[lbl, above=1mm, midway, text width=24mm]
    {HTTP + token\\office LAN} (server.west);
  \draw[flow] (server.east) -- (core.north west);
  \draw[flow] (cli.south east) -- (core.north);
  \draw[flow] (core) -- (lance);
  \draw[flow] (md) -- (lance);
  \draw[flow] (raw) -- (md);
  \draw[flow] (core.east) -- node[lbl, above=1mm, midway, text width=26mm]
    {HTTP, forwarded\\to the pool} (ollama.west);
  \draw[flow, <->] (ollama.east) -- (chat.west);
  \draw[flow, <->] (ollama.east) -- (embed.west);
  \end{tikzpicture}}
  \caption{Where everything runs. Only questions and answers cross the network;
           the index stays beside the documents.}
  \label{fig:corpus-system}
\end{figure}

Two consequences are worth holding onto. First, \textbf{the corpus is not the
model}: the laptop's SSH tunnel of Section~\ref{sec:remote-access} reaches the
model on port 11434, but the corpus is a separate service on port 8765 and
needs its own path. Second, \textbf{retrieval happens where the data is}. Only
a question goes over the wire and only an answer comes back --- kilobytes, not
the gigabytes it would take to replicate the index.

\subsection{Ingest: documents in, index out}
Ingest runs once per document, and is cached. Parsing is the slow step, so its
output is stored in \cmd{md/} keyed on the source file's sha256 --- re-running
\cmd{./rag ingest} over a growing folder only touches files whose bytes
changed.

\begin{figure}[!ht]
  \centering
  \begin{tikzpicture}[node distance=6mm,
    blk/.style={draw=atBorderStrong, rounded corners=2pt, fill=atBg,
                align=center, inner sep=4pt, font=\scriptsize,
                minimum height=8mm},
    store/.style={blk, fill=atBgSubtle},
    comp/.style={blk, fill=atBgMuted},
    model/.style={blk, fill=atAccent!30},
    flow/.style={-{Stealth[length=2mm]}, draw=atFgMuted, thick},
    lbl/.style={font=\tiny, text=atFgMuted, align=center},
  ]
  \node[store] (src) {\cmd{\char`~/work/raw}\\PDF DOCX XLSX\\CSV MD TXT};
  \node[comp, right=of src] (disc) {discover\\skips venvs,\\\cmd{node\_modules}, \cmd{.git}};
  \node[comp, right=of disc] (parse) {parse\\to markdown};
  \node[comp, right=of parse] (fix) {repair\\ligatures, tidy};
  \node[store, right=of fix] (cache) {\cmd{md/}\\cache};

  \node[comp, below=14mm of cache] (chunk) {chunk\\1000 tok, 150 overlap\\heading-aware};
  \node[model, left=of chunk] (emb) {embed\\\cmd{bge-m3}};
  \node[store, left=of emb] (tbl) {LanceDB\\\cmd{chunks} table};
  \node[comp, left=of tbl] (idx) {build indexes\\BM25 always\\ANN $\geq$ 2048 rows};

  \draw[flow] (src) -- (disc);
  \draw[flow] (disc) -- (parse);
  \draw[flow] (parse) -- (fix);
  \draw[flow] (fix) -- (cache);
  \draw[flow] (cache.south) -- (chunk.north);
  \draw[flow] (chunk) -- (emb);
  \draw[flow] (emb) -- (tbl);
  \draw[flow] (tbl) -- (idx);

  \node[lbl, below=1mm of parse, text width=26mm] {scans flagged,\\not indexed};
  \node[lbl, below=1mm of chunk, text width=30mm]
    {title + heading path\\prefixed to each chunk};
  \end{tikzpicture}
  \caption{Ingest. Parsing is cached on the source's sha256, so re-running only
           touches documents whose bytes changed.}
  \label{fig:corpus-ingest}
\end{figure}

\subsubsection*{Why the chunks look the way they do}
Every chunk is prefixed with its document title and heading path before it is
embedded. A passage reading ``the limit is 40\,$^{\circ}$C'' is useless on its
own; ``TB 754 $>$ 6.~Harmonic limits $>$ 6.11 --- the limit is
40\,$^{\circ}$C'' both retrieves and answers correctly. The prefix is what gets
embedded and BM25-indexed; the bare text is what the model is shown.

Tables are kept whole where possible, and split on row boundaries with the
header repeated when they are too large --- a spreadsheet row severed from its
column headings cannot be read by anything.

\subsection{Query: three model calls, not one}
A question triggers an embedding call, then a rerank call, then the answer
call. The database and the model never talk to each other; Python moves text
between them.

\begin{figure}[!ht]
  \centering
  \resizebox{\linewidth}{!}{%
  \begin{tikzpicture}[
    blk/.style={draw=atBorderStrong, rounded corners=2pt, fill=atBg,
                align=center, inner sep=4pt, font=\scriptsize,
                minimum height=8mm},
    comp/.style={blk, fill=atBgMuted},
    model/.style={blk, fill=atAccent!30},
    zone/.style={draw=atBorderStrong, dashed, rounded corners=4pt, inner sep=6mm},
    flow/.style={-{Stealth[length=2mm]}, draw=atFgMuted, thick},
    lbl/.style={font=\tiny, text=atFgMuted, align=center},
  ]
  % retrieval
  \node[blk]   (q)    at (0,0)      {question};
  \node[model] (qe)   at (2.5,0)    {\textbf{1.} embed\\\cmd{bge-m3}};
  \node[comp]  (vec)  at (5.6,1.15) {vector search\\cosine, top 40};
  \node[comp]  (bm)   at (5.6,-1.15){BM25 search\\top 40};
  \node[comp]  (rrf)  at (9.1,0)    {fuse\\reciprocal rank};
  \node[comp]  (prom) at (12.1,0)   {promote a\\dominant\\BM25 hit};
  \node[zone, fit=(vec)(bm),
        label={[font=\tiny, text=atFgMuted]below:LanceDB --- one table, two indexes}] (lz) {};

  % generation
  \node[model] (rr)  at (12.1,-4.3) {\textbf{2.} rerank\\\cmd{gpt-oss:120b}\\listwise, 1 call};
  \node[model] (ans) at (8.0,-4.3)  {\textbf{3.} answer\\\cmd{gpt-oss:120b}};
  \node[blk]   (out) at (4.2,-4.3)  {cited answer\\+ sources};

  \draw[flow] (q) -- (qe);
  \draw[flow] (qe.east) to[out=40,in=180]  (vec.west);
  \draw[flow] (qe.east) to[out=-40,in=180] (bm.west);
  \draw[flow] (vec.east) to[out=0,in=130]  (rrf.north);
  \draw[flow] (bm.east)  to[out=0,in=-130] (rrf.south);
  \draw[flow] (rrf) -- (prom);
  \draw[flow] (prom) -- node[lbl, right=1mm, midway, text width=20mm]{40\\candidates} (rr);
  \draw[flow] (rr) -- node[lbl, above=1mm, midway]{top 8 passages} (ans);
  \draw[flow] (ans) -- (out);

  \node[lbl, text width=40mm] at (5.6,2.5) {embeddings blur exact identifiers};
  \node[lbl, text width=30mm] at (1.5,-1.9) {BM25 catches ``754'', ``IEC 60287''};
  \node[lbl, below=2mm of ans, text width=38mm]
    {\cmd{-{}-fast} skips step 2:\\$\approx$7\,s instead of $\approx$20\,s};
  \end{tikzpicture}}
  \caption{A query. Three model calls: embed the question, rerank the
           shortlist, then answer from the survivors.}
  \label{fig:corpus-query}
\end{figure}

\subsubsection*{Which model does what}
There are exactly four model calls in the whole system, and \textbf{ragkit}
makes all of them. The two models are not interchangeable: \cmd{bge-m3} only
turns text into vectors --- it cannot generate --- so anything that produces
words is \cmd{gpt-oss:120b}.

\begin{table}[!ht]
  \centering
  \renewcommand{\arraystretch}{1.25}
  \begin{tabularx}{\linewidth}{@{}l l l X@{}}
    \toprule
    \textbf{Call site} & \textbf{Wrapper} & \textbf{Model} & \textbf{When} \\
    \midrule
    \cmd{rag.py:87}       & \cmd{llm.embed} & \cmd{bge-m3}       & ingest: embed every chunk \\
    \cmd{retrieve.py:62}  & \cmd{llm.embed} & \cmd{bge-m3}       & query: embed the question \\
    \cmd{retrieve.py:192} & \cmd{llm.chat}  & \cmd{gpt-oss:120b} & rerank the shortlist \\
    \cmd{answer.py:36}    & \cmd{llm.chat}  & \cmd{gpt-oss:120b} & write the cited answer \\
    \bottomrule
  \end{tabularx}
  \caption{The four model calls, and which model serves each.}
  \label{tab:corpus-model-calls}
\end{table}

\cmd{llm.embed} posts to \cmd{/api/embed}, \cmd{llm.chat} to
\cmd{/v1/chat/completions} --- same Ollama, same connection, different model
name in the payload. That is why Figure~\ref{fig:corpus-system} shows one arrow
into Ollama and Ollama reaching both models.

Where the LanceDB search sits in that list is: nowhere. The lookup itself
(vector + BM25) uses no model at all. It is the steps on either side of it that
do --- embedding the question before, reranking after. Likewise ``the database
and the model never talk to each other'' means exactly that: LanceDB and
\cmd{gpt-oss:120b} have no link, and \textbf{ragkit} is the thing in the middle
that talks to both.

\subsubsection*{How the model actually receives the passages}
As plain text in the prompt. There is no special channel and nothing is loaded
into the model. The retrieved passages are concatenated into one numbered block
and sent as an ordinary chat message to \cmd{/v1/chat/completions}:

\begin{center}
\begin{tikzpicture}
\node[draw=atBorderStrong, rounded corners=2pt, fill=atBgSubtle,
      text width=110mm, align=left, inner sep=6pt, font=\scriptsize] (p) {
\textbf{system}\ \ Use only the numbered sources. Cite as [3]. If the sources
do not contain the answer, say so.\\[3pt]
\textbf{user}\ \ \ \ \ Sources:\\
\phantom{xx}[1] 0754 AC side harmonics $>$ 6.~Harmonic limits $>$ 6.11.2 \dots\\
\phantom{xx}distortion factor of the centred interharmonic subgroup \dots\\
\phantom{xx}[2] \dots\\[2pt]
\phantom{xx}Question: What harmonic order range does brochure 754 consider?};
\end{tikzpicture}
\end{center}

The model sees one flat sequence of tokens. It has no filesystem access, no
connection to LanceDB, and no memory of previous questions --- each question
builds a fresh prompt. Three things follow:

\begin{itemize}\itemsep1pt
  \item \textbf{Citations cannot drift.} The model copies back the \cmd{[n]}
        label it was given; the mapping from \cmd{[n]} to a file path is done
        in Python afterwards.
  \item \textbf{Refusal is honest.} ``The sources do not contain any entry for
        brochure 9999'' is an accurate report about its input.
  \item \textbf{Retrieval is the ceiling.} If a passage does not reach the
        prompt, the model cannot know it, however capable it is.
\end{itemize}

\subsubsection*{Why context size is not the constraint}

\begin{table}[!ht]
  \centering
  \renewcommand{\arraystretch}{1.25}
  \begin{tabularx}{\linewidth}{@{}X r@{}}
    \toprule
    & \textbf{Tokens} \\
    \midrule
    System message                        & 128 \\
    Retrieved passages (4) + question     & 4\,103 \\
    Typical query at \cmd{k=8}            & $\approx$8\,000 \\
    \midrule
    One brochure in full (TB 754)         & 74\,550 \\
    Model context window                  & 131\,072 \\
    \bottomrule
  \end{tabularx}
  \caption{Token budget of a query against the context window.}
  \label{tab:corpus-context}
\end{table}

A default query uses about 6\,\% of the window. Stuffing the context instead
would fit one brochure and a half out of roughly a thousand. That is the whole
argument for retrieval: the problem was never how much text fits, it is picking
the right four thousand tokens.
